% ============================================================
%  Harald Høffding, POSITIVISM AND MATERIALISM [Positivisme og
%  Materialisme. Med særligt Hensyn til Maudsleys Foredrag ...].
%  Nyt dansk Maanedsskrift, ed. Vilhelm Møller, vol. 2 (April--September
%  1871), pp. 193--208 (sections 1--2) and 309--333 (sections 3--4).
%  TRANSLATION (English)
%
%  Translated from transcription.tex (Danish) in this directory, whose
%  header records how its text was established from the Internet Archive
%  scan (every page transcribed at the image, diffed against the text
%  layer, and read a second time at the image without it).  Method:
%  ../../../TRANSLATION-PLAYBOOK.md.  No published translation was
%  consulted (none is known).
%
%  ---- REGISTER ----
%  Scholarly, moderately literal, readable English; American spelling.
%  Long periods are broken only where English will not carry them, never
%  across a page marker or a footnote anchor.
%
%  ---- TERMINOLOGY ----
%  Aligned with this repo's translations of Høffding's Philosophie og
%  Theologie (1866) and Filosofien og Darwinismen (1874):
%    Sjæl, sjælelig = mind, mental (Maudsley's "mind and body"); "soul"
%      only where the sense demands it: p. 200 „vegetativ Sjæl“; pp.
%      311--312 the spiritualist doctrine of the immaterial soul, Plato,
%      Lemoine; p. 319 „besjælet“ = ensouled (% TR: at each site);
%      Sjæleliv = mental life;
%    Legeme, legemlig = body, bodily; Sindssygdom = mental illness;
%      Afsindighed = insanity; Nervelidelse = nervous disorder;
%    Aand = spirit; Heelhed = the whole; Betingelse = condition;
%    Videnskab = science; Erkjendelse = cognition; Erkjendelseslære =
%      theory of cognition; Forsken = inquiry; Anlæg = disposition;
%    Virkelighed / Realitet = actuality / reality; Væsen = essence;
%    Forestilling = representation; Kraft / Stof = force / matter;
%    Slægten = the race (the human lineage, in the discussion of
%      heredity); Tilregnelighed = accountability.
%
%  ---- CONVENTIONS ----
%  - Page markers \opage{N} and "% --- p. N ---" at the same joints as in
%    transcription.tex.
%  - The same heads, subtitle, sections 1.--4., rules, footnotes (given
%    whole at the mark, including those that run over onto the next
%    page), \emph spans (letterspacing in the original), \textit spans
%    (antiqua), epigraph, signatures and closing rules as the Danish.
%  - Quotations in other languages stay as printed: French, Latin and
%    English in \textit, French in its roman guillemets »...« (with the
%    print's wrong-way closings); German as printed, in roman (the print
%    sets it in Fraktur).  The Mill epigraph (p. 309) is English in the
%    original.  English authors whom Høffding quotes in Danish (Maudsley,
%    Mill, Lewes) are translated from his Danish, not restored from the
%    English originals.  Book and journal titles stay in their language.
%  - Quotation marks „...“ -> ``...''.  „(Fortsættes.)“ = (To be continued.)
%  - Translator's decisions are recorded in "% TR:" comments at the site.
%
%  ---- DEFECTS: DANISH AS PRINTED, ENGLISH CORRECT ----
%  The printer's errors listed in the transcription's header are rendered
%  in their evident sense and logged in % TR: at each site (p. 204
%  „Sammensat“; p. 314 „ellér“; p. 318 „tidligere“; p. 327 „theoride“ =
%  Maudsley's ``theroid''; p. 329 ‚Arvelighed“).
%
%  ---- HOW THE TRANSLATION WAS CHECKED (2026-09-27) ----
%  (1) Five translators (pp. 193--200, 201--208, 309--316, 317--324,
%  325--333), each re-reading its English against the Danish sentence by
%  sentence.  (2) Mechanical audit (tr_audit.py): per printed page the
%  same \opage joints, \emph and \textit spans, footnotes, section
%  numbers, centre blocks and paragraph breaks as transcription.tex, and
%  an English/Danish word ratio within 0.95--1.60.
%  (3) Independent review: five readers who translated none of it read
%  every sentence of the Danish against the English and proposed 13
%  corrections, all applied.  Sense: p. 205 „opfattede“ = interpreted, not
%  took up; p. 316 a dropped concessive „jo“; p. 320 „ganske vist ...
%  men“ concessive (to be sure); p. 321 „maa vist betvivles“ = must surely
%  be doubted, and „hiin“ = the current (not "the former"); pp. 326, 328
%  „al videre moralsk Følelse“ = any real moral feeling (not "all
%  further").  Terminology: Sjæl = soul at four sites on pp. 311--312
%  (above), while the reviewers judged "mind" right in Høffding's own
%  argument on pp. 313--316 and in his report of Maudsley on p. 310;
%  p. 328 „Samlen og Sammenfatten“ = gathering and comprehending, echoing
%  pp. 313, 316.  English: p. 329 a misbuilt cleft.
% ============================================================
\documentclass[12pt,a4paper,openany]{book}
\usepackage[utf8]{inputenc}
\usepackage[T1]{fontenc}
\usepackage{libertinus}
\usepackage{libertinust1math}
\usepackage{textalpha}
\usepackage{graphicx}
\usepackage[english]{babel}
\usepackage[protrusion=true,expansion=true]{microtype}
\usepackage{setspace}
\usepackage[margin=1.2in]{geometry}
\usepackage{enumitem}
\usepackage{fancyhdr}
\setlength{\headheight}{28pt}
\addtolength{\topmargin}{-13.5pt}

% Footnotes as the printing sets them: *).
\renewcommand{\thefootnote}{*)}

% The article's numbered sections: the number alone, centred on its line,
% exactly as in transcription.tex.
\newcommand{\secnum}[1]{\par\begin{center}#1\end{center}\par\nobreak}

\usepackage{hyperref}
\hypersetup{
  pdftitle={Positivism and Materialism},
  pdfauthor={Harald Høffding},
  hidelinks
}

\pagestyle{fancy}
\fancyhf{}
\fancyhead[LE,RO]{\thepage}
\fancyhead[RE]{\textit{Harald Høffding: Positivism and Materialism}}
\fancyhead[LO]{\textit{Nyt dansk Maanedsskrift, vol.\ 2}}

\setstretch{1.3}

\title{\textbf{Positivism and Materialism}\\[8pt]
  \large [Positivisme og Materialisme]}
\author{Harald Høffding}
\date{\textit{Nyt dansk Maanedsskrift}, ed.\ Vilhelm Møller, vol.~2
  (April--September 1871), pp.~193--208, 309--333\\[10pt]
  \small Translated from the Danish}

% ---- original-page markers ----
% \opage{N} marks where printed page N begins, at the same joint as in
% transcription.tex (mid-sentence where the page turns mid-sentence).
\usepackage{marginnote}
\newcommand{\opage}[1]{\marginnote{\footnotesize\textit{[#1]}}}
\newcommand{\apage}[1]{\marginnote{\footnotesize\textit{[#1]}}}

\begin{document}
\maketitle
\thispagestyle{empty}
\clearpage

% --- p. 193 ---
\opage{193}\begin{center}\textbf{\large Positivism and Materialism.}\end{center}

% footnote *) on p. 193
{\small With Particular Reference to \emph{Maudsley's} Lecture: ``On the Relation between Mind and Body and between Mental Illnesses and Other Disorders of the Nervous System''. (Translated in ``Ugeskrift for Læger''. Sept.--Oct. 1870)\footnote{Later published as a separate offprint.}.\par}

\begin{center}\rule{1.8cm}{0.4pt}\end{center}

\secnum{1.}

It is only in the modern age that one can speak of sciences.
Antiquity and the Middle Ages really knew only one,
all-encompassing science. The reason for this lies in the essentially
% TR: „det væsenlig forskjellige Synspunkt“ is singular in the Danish; rendered plural, since the next sentence names „Det ene Hovedsynspunkt“ (one of two)
different standpoints from which existence can be conceived.
The one chief standpoint is that of the whole, the one, final end
in which everything gathers together for comprehending thought, and
which must therefore be the concluding point and summit of every
worldview. Antiquity and the Middle Ages adopt \emph{this}
standpoint alone, and therefore for them science is only
one. From the idea of the whole, which in antiquity is essentially seen as
the idea of the beautiful, in the Middle Ages as the idea of the good,
one then seeks to derive the whole system of conditions that
forms the foundation for the actualization of the idea; the conditions
are seen as contained in the idea itself, without any felt need for
a distinct inquiry with distinct laws and methods in order to
investigate them. In the modern age, by contrast, thought has
turned especially toward the conditions as such; it has
% --- p. 194 ---
\opage{194}seen the impossibility of arriving at the cognition of them except
through a peculiar inquiry occupied specifically with them,
and little by little, therefore, one circle of conditions after
another detached itself from the great whole as the object of a
single, self-contained science with its own method and its own principles.
Experience is installed in its rights and brought under
exact laws.

This dissolution of the abstract unity of science
could not take place without resistance on the part of those
who believed that with this absolute unity they would lose all scientific
unity in cognition. At the beginning of the modern
age, to be sure, it was chiefly the medieval Aristotelians who opposed
empirical science; and for them the ideal of unity had
stiffened into a faith in authority that let them set
% footnote *) on p. 194
% TR: in the footnote the bibliographic data „2te Ausg. S. 213“ are rendered in English (2nd ed., p. 213); the titles stay as printed; „Galilæi“ kept as printed
Aristotelian pronouncements above the testimony of their own eyes\footnote{\emph{\textit{Galilæi}} \textit{(Dialogus de systemate mundi} cited by \emph{Feuerbach}: Pierre Bayle. 2nd ed., p. 213) relates that when a famous physician in Venice had demonstrated in a lecture that the largest nerve trunk issued from the brain, he asked an Aristotelian who was present whether he still believed that the nerves originated in the heart and not in the brain; the latter then answered that he would have declared himself convinced by what he now saw before his eyes, if Aristotle had not expressly taught the opposite: \textit{Eqvidem ita aperte rem oculis subjecisti, ut, nisi textus Aristotelicus aperte nervos ex corde deducens obstaret, in sententiam tuam protracturus me fuisses!}}.
But at the beginning of this century the speculative idealism in Germany attempted
once again to develop everything from a single
% TR: „Baco“ kept as printed (Francis Bacon)
starting point, the idea. Schelling declared that Baco had
corrupted philosophy, and Newton physics! Yet nothing has
been able to halt the victorious advance of empirical science;
step by step it has pressed forward and has encompassed everything
in the net of the exact method.

This great power of the spirit has not escaped the danger
that, as the ancients so profoundly teach us, threatens every
fortunate and victorious power: that of overstepping its limit and its
% --- p. 195 ---
\opage{195}justification. But it must be stressed that the limit here
is not easy to draw. In this domain there can be no question
of outward, finite limits. The limit must lie
originally in the very principle from which one proceeds. Modern
empirical science does not ask about the ideal
significance and worth of its objects; it seeks only to determine their
elements and inner connection; it aims at a
\emph{how} and not at a \emph{why}. But then it follows
of itself that the result, too, can be only a cognition
of the elements and their mutual connection, --- that, therefore,
a total, concluding cognition of existence requires
as its supplement a cognition that holds fast to the
ideal standpoint and proceeds from the total unity. Only
the interaction and union of these standpoints
can lead to the goal; they are the two ends that must be joined
for the circle to be closed. Empirical science therefore oversteps
its limit when it believes straightway that it has
solved the \emph{whole} problem of existence. It holds only one
of the threads in its hand.

At this point lies the germ of materialism.
The idea, the total unity, is the determination of form, over against
% TR: „Stoffets, Materiens Bestemmelse“ -- Stof and Materie are both "matter"; rendered "of matter, of the material"
which the system of conditions represents the determination of matter, of the
material, and the peculiarity of materialism
now consists in overlooking the independence of form, of the idea, and
deriving it solely from the connection of matter. It regards
matter, then, not merely as the condition of form, but
as the cause of form. To the uncritical eye the materialistic tendency
thus lies so close to empirical science
that it is no wonder that in our day it dares to regard itself
as the latter's only legitimate daughter. Among natural scientists who
do not have a sharp and sure eye for the essence
of their science, materialism readily makes its appearance as soon as they
seek to raise themselves from the detail of special investigations to a more
comprehensive total view (as a rule in a very popular form);
they forget that for such a view other elements come
into consideration than those that the single science can encompass.
% --- p. 196 ---
\opage{196}This belongs among the unfortunate consequences of the division of labor in
the domain of cognition.

Yet it is not only where the final,
concluding total view is concerned that the materialistic tendency,
as we have characterized it, can assert itself.
Within empirical science itself, too, the same relation
between condition and cause, matter
and form, repeats itself; \emph{there}, too, then, lies a possibility of opposite
directions. The more the great advances of modern science
in every domain are owed to the exact conception of the conditions,
the nearer lies the misunderstanding that the thing
itself is nothing but the conditions. Materialism is
the reaction against the dualism of earlier times, a protest against
the external separation of the essence from its conditions.
What it combats is really only the opinion that there
should be hidden, as it were, behind the phenomena a certain
mystical force that was independent of the conditions. It
therefore strives to lead everything back to the elementary and to let
it arise from the composition of the elementary. The higher, more
concrete science must then constantly struggle against encroachments
from the lower, more abstract and elementary one. In physics
that direction is termed materialistic which denies the peculiarity of the individual
mechanical forces and regards them only
as special forms of elementary motion. If
this is extended further, not only will chemistry be dissolved into mechanical physics,
but physiology too into chemistry, psychology into physiology.
The concrete and peculiar is at every point to be
seen only as the product of the abstract, general conditions, the
individual unity only as a composition of mechanically
acting elements.

We are speaking here only of the materialistic \emph{tendency}.
Materialism cannot be carried through completely. Pure
matter would be one with pure nothing, from which nothing can
come. The essence of materialism lies rather in the relative,
as indeed the determination of matter is always relative, since
what is matter in one respect is form in another.
% --- p. 197 ---
\opage{197}The justification of relative materialism lies, according to the
foregoing, in expressing the truth that matter explains
form \emph{insofar as} the latter rests on and depends on the former.
But when the tendency goes further, the self-contradiction comes
plainly to light.

For form, the idea of the whole, does not come in from outside,
but shows itself within the very world of
conditions. It finds its particular expression --- although it is
present everywhere --- where a peculiar circle of laws and
phenomena unfolds on the basis of more abstract and
elementary presuppositions, and where therefore the domain of a new and higher
science begins. These different circles
form, as it were, a nodal line, where the quantitative transitions
condense into qualitative centers, from which the movement,
through continual quantifying, constantly passes on to
% footnote *) on p. 197, continued at the foot of p. 198; given whole here at its mark
% TR: in the footnote Lewes is quoted in Danish; translated from the Danish, the English original not reconstructed; „ɔ:“ (id est) rendered "i.e."
other and yet other nodal points\footnote{Compare with this what \emph{Lewes} says in his biography of Goethe about the latter's theory of colors. It is mistaken as a contribution to optics, because it disregards the exact quantitative method. ``Without the free control of the scales the chemical experiment can never become quantitative, and without quantitative knowledge there can be no chemical science, strictly speaking, but only qualitative, i.e.\ approximate, knowledge. No degree of observation can make observation exact when it cannot be measured. Though one should observe falling bodies for all eternity, without mathematics one will never arrive at the law of gravitation.'' But on the other hand Goethe's theory of colors is very valuable when one looks only to the conception of colors as such, without regard to their mechanical conditions. ``For artists and physiologists, that is to say for those who have an interest chiefly in the phenomena of color as sense impressions, and who demand qualitative rather than quantitative knowledge, his works have great value''. --- On this it is to be remarked that it is only from an abstract quantitative standpoint that the qualitative can be made one with the approximate. According to Lewes's own characterization, by contrast, the qualitative designates the phenomenon as a whole and the significance, the worth, that it has as such, and this becomes here, where colors are in question, one with the ideal, psychological and aesthetic significance. Naturally the different standpoints cannot be kept absolutely apart from one another; it is precisely this that is designated by the expression nodal line.}. Materialism
% --- p. 198 ---
\opage{198}now grounds itself here on an optical illusion, inasmuch as,
because the qualitative can be determined through quantifying,
it believes that the qualitative in itself rests only on differences of magnitude.
But since those nodal points actually exist in
existence, the materialistic tendency constantly comes to
contradict its own principle, comes to smuggle in qualitative
determinations while it believes it is explaining everything quantitatively.

Viewed theoretically, then, materialism in our time is easy
to explain; it is bound up with what, viewed from another
side, constitutes the greatness and strength of our time. It is a
phenomenon of the time, whose essence and origin one need only
make clear to oneself in order to be able to overcome it. In practice
materialism no doubt appears with greater tenacity and strength;
yet here too the same must hold. From the theological
side materialism has been summarily condemned as heresy,
inasmuch as it is regarded not as an error in cognition,
but as a doctrine formed by the evil will for its own
% footnote *) on p. 198, continued at the foot of p. 199; given whole here at its mark
% TR: in the footnote the German quotations are kept as printed (in roman, as the Danish sets them), including „mekanischen“ and „Rechtsbewußtseins“; „2te Aufl. S. 216“ and Høffding's „(S. 27)“ rendered "2nd ed., p. 216" and "(p. 27)"; Fosterfordrivelse = abortion
extenuation\footnote{Cf. ``Briefe gegen den Materialismus von \textit{Dr.} F. Fabri, Missionsinspektor''. 2nd ed., p. 216: ``Die Lehre der Materialisten von einer mekanischen Denknothwendigkeit ist im Grunde nichts als das Feigenblatt ihrer die moralische Verantwortlichkeit läugnenden Unfreiheit des Willens. Der verkehrte Wille vollendet sich in der verkehrten Theorie''. The author holds, namely, that all knowing proceeds from an immediate faith, which in turn is determined by man's will. He thereby really makes his whole book superfluous, while at the same time it becomes consistent when he attributes unethical motives to his opponents, which he drives to a fanatical extreme. Thus he mentions in a note that a ``Verein badischer Aerzte zur Förderung der Staatsarzneikunde'' had named Büchner, author of the well-known materialistic work ``Kraft und Stoff'', its honorary member. It would not, he adds (p. 27), be uninteresting to learn something about the motives that guided the society for ``Staatsarzneikunde'' in that step: ``Doch wir erinnern uns so eben, daß der Verfasser von ``Kraft und Stoff'' die Abtreibung der Leibesfrucht als ein Recht der Mutter wieder vindicirt hat, und daß er die Freiheit des Willens und also die sittliche Zurechnungsfähigkeit mit Herrn Vogt läugnet! Es kann aber nicht in Abrede gestellt werden, daß dies Behauptungen sind von solcher Tragweite, die allerdings nicht nur die gesammte gerichtliche Medicin, sondern unser ganzes Rechtsbewußtseins und alle sociale Ordnungen in Grund und Boden reformiren müßten. Und bei der Lage des Falls bleibt wenigstens bis zur Führung des Gegenbeweises (!) nichts übrig, als anzunehmen, jener Carlsruher Verein habe in Herrn Büchner den verkannten Apostel einer solchen schönen Zukunft ehren wollen''. The materialists did not handle their opponents gently, but one would search their writings in vain for such venomous insinuations as the one cited here. I hardly need remark that Büchner by no means defends abortion; he merely mentions, in order to show how moral representations vary at different times, that while we regard it as a crime, the Romans found nothing immoral in it. It is Fabri himself who, in his zeal for the ``mission'' he has taken upon himself against materialism, finds in this an ``assertion of abortion as a right of the mother''!}. For philosophical criticism, on the other hand,
% --- p. 199 ---
\opage{199}the practical corrective stands in close connection with the
theoretical and does not appear in isolation. Thereby criticism
also becomes able to acknowledge what is justified in the materialistic
tendency.

In our day a grand attempt has been made to
lay what is justified in the materialistic tendency at the foundation
of a whole worldview. So-called positivism,
whose founder is \emph{Auguste Comte}, and which has won
wide adherence among the French and English philosophers and
natural scientists, aims to carry through a philosophy grounded
immediately on the principle of empirical science, ---
a ``positive'' philosophy, because it is supported by the positive
sciences. A consideration of this direction will demonstrate more clearly
than general reflections the peculiar character of relative materialism.
% TR: „Maudsleys interessante Foredrag, der ... have vakt ... høre“ -- plural verbs, so "lectures" here (the subtitle's fixed wording keeps "Lecture")
\emph{Maudsley's} interesting lectures,
which have rightly attracted attention in a wide
circle, belong, as regards the whole underlying train of thought,
to the general positivist direction.
% --- p. 200 ---
\opage{200}\secnum{2.}

Auguste Comte sets his positive philosophy in opposition
to the theological and metaphysical philosophizing of earlier times.
At the first, imperfect stage of consciousness man sees
everything in nature as proceeding from a cause similar to the one that
brings about his own movements and actions, that is, from a
thought and will lying behind the phenomena. Since he
does not yet have any clear consciousness of necessary lawfulness,
he sees the will at work in nature as an arbitrary
will, which does with things just as it pleases; and since
this in turn presupposes a higher power than the human,
which has no absolute dominion over nature,
that arbitrary will becomes for man the divine omnipotence,
to which he relates himself through religion.
But through closer, sustained occupation with the phenomena
there gradually arises the cognition that it is not
absolute arbitrariness but fixed order and regularity that reign
in their course. In place of the almighty, personal
will, the inner essence of things then appears to consciousness
as the real cause. Behind the phenomena one imagines
hidden forces, in which their final explanation
is sought; each of these forces rules in its circle of phenomena,
and the system of all the forces forms an inner
world which produces and sustains the outer, visible
world. In place of that world of gods in which the
immediate consciousness sought the explanation of phenomena,
there is here, then, set a world of metaphysical beings \textit{(entités métaphysiques)}.
It is no longer dryads that animate
the trees and guide their growth and life, but the cause of
% TR: „vegetativ Sjæl“ -- here only "soul" will do (the vegetative soul of the Aristotelian tradition)
plant as well as of animal life is sought in a ``vegetative soul''
or a ``life force''. The analogy with human consciousness
recedes further here. Yet it is common
to the religious and the metaphysical stage of consciousness that they
both explain the outer by the inner, that they make man's
essence the basis for the explanation of the world, instead of
conversely explaining man by the world.

% --- p. 201 ---
\opage{201}»\textit{Le véritable esprit général de toute philosophie théologique
ou métaphysique consiste à prendre pour principe,
dans l'explication des phénomènes du monde extérieur,
notre sentiment immédiat des phénomènes humains; tandis
que, au contraire, la philosophie positive est toujours caracterisé,
non moins profondément, par la subordination
nécessaire et rationelle de la conception de l'homme à celle
du monde}« \textit{(Auguste Comte: Cours de philosophie positive.
2éme éd. Tome III. p. 188)}.

The third period of consciousness is marked by the dominion of the
positive, scientific method. Only because the religious
and the metaphysical consciousness proceeded from the essence of
man could they suppose that they attained a cognition of the essence
and causes of things, an insight into the origin and
destination of the world. Thought renounces this once it has
arrived at »\textit{l'état positif}«. It now keeps to the positive,
that is to say, to the factually given: »\textit{positif, c'est à dire
fondé sur des faits bien constatés}». But the factually given
is only the phenomena. The sole task of thought is then to
find their actual laws, that is to say, the invariable relations
in which they succeed one another and resemble one another.
It is, on the other hand, impossible to know their nature and essence in
and for themselves. The essence in and for itself would be the unconditioned
and absolute; but the positive is always conditioned and relative.
Positive philosophy must therefore set up as its absolute
fundamental principle that we can cognize nothing absolute.

What Comte calls »\textit{la loi des trois états}« thus leads us
into the peculiar essence of positivism. But closer determinations
prove to be necessary. We cognize only the positive,
i.e.\ (as Comte determines it) the relative, the conditioned.
But does this mean that only the relative is, that no
absolute principle \emph{exists}, or does it mean only that we cannot
\emph{cognize} it? When everything we cognize is conditioned, is
that to say that everything is only conditions, so that the conditioned
is only the derived property of the conditions? In other
% --- p. 202 ---
\opage{202}words, there lies in the principle of positivism a possibility of skepticism
as well as of materialism.

Since Comte rejects all cognition of inner causes,
of the essence and principle of things, his endeavor aims only
at discovering the laws of phenomena, i.e.\ the conditions
under which they arise. These laws or ``general
facts'' \textit{(faits généraux)}, which hold for a particular
circle of phenomena, point back in turn to more
general relations, of which they are only a particular form,
and by thus always going back one comes at last
from the concrete, intricate and composite phenomena to the
simplest, most elementary and therefore most general laws.
The mathematical laws embrace all spheres of existence,
because they presuppose solely the pure relations of magnitude. »\textit{La loi
des trois états}« is here found completely realized: mathematics
has long been a positive science. This does not, however,
hold of the other sciences. Here particular principles are
still always presupposed which are not reduced
to simple positive phenomena or relations between phenomena.
The completion of positivism \emph{then} depends on this
happening: that every peculiarity, every particular principle,
be traced back to elementary presuppositions, and that these in turn
be seen to proceed from the preceding, more abstract and therefore
more comprehensive circle of laws. Each such circle, to be sure,
always presupposes a method of experience of its own; but it
cannot become systematic except by being derived from the more
elementary circles.

What is indicated here goes in actuality beyond
what is contained in the original program of positivism.
Positivism must necessarily cease where the ``positive''
ends, and this happens as soon as the factual connection of
the phenomena ends. The law Comte believes he has
found, that the most general properties, which hold in the
widest extension, are at the same time the most elementary and abstract, ---
a law which, as has rightly been remarked, is only the
well-known logical proposition of the inverse relation between
% --- p. 203 ---
\opage{203}the content of a concept and its extension\footnote{The content of a concept is the marks it specifies; its extension is everything to which these marks apply and which therefore falls under the concept. Now the greater the content (the more marks the concept specifies), the smaller the extension (the individuals that possess the concept's marks and so can be brought under it). E.g.\ the concept animal --- the concept horse: the latter contains more marks, but embraces fewer individuals than the former.} --- this law contains
no guarantee that the real (``positive'') explanation
of phenomena could be found by always going back
from the concrete and combined to the abstract and
elementary. Comte has here, without further ado, transferred a logical
relation to the real and positive, and has thereby placed himself
at a standpoint similar to that of speculative philosophy, which
derives the real from the logical. He here explains, to
use his own words, the world by man, not man
by the world.

The task he has set himself is to develop the
fundamental concepts on which all our thinking and cognition
rest, and he declares expressly that these fundamental
concepts are not themselves grounded on any ``preceding
intellectual system''. But in the use he makes of
that logical law he has an unconscious presupposition which is not
drawn into the investigation. Here positivism then shows itself
as dogmatism, and thereby passes at the same time over into materialism.
If one has granted that the logical proposition of
the inverse relation between the content and extension of a concept
can be applied immediately in the explanation of nature,
then one has thereby granted the materialistic principle itself.
But in this immediate application it is overlooked that
the proposition undergoes a considerable change by being transferred
in this way to the domain of nature. Logically seen,
the abstract concept is the higher, insofar as it includes
the more concrete concepts within it; but if one does not remain
standing at a purely formal mode of consideration, the
% --- p. 204 ---
\opage{204}higher concept must be developed in such a way that it does not merely arise
by subtraction of the individual characteristic marks, but contains
these within itself as possibilities, as moments, so that it
gives the final, concluding explanation of the concrete, over against
which it has the advantage of not being bound to the
individual and finite, and yet, ideally seen, of including this within
itself. The return from the higher concept, now not more abstract but
more concrete, to the lower, more limited one does not then take place
by a mere addition of the individual specific characteristic marks;
rather, these are developed from the former as particular, relative
expressions of it. As Comte conceives the higher concept
and applies it in the explanation of phenomena, it is
the merely abstract general concept, emptied by the subtraction of the particular.
He thus regards the organic merely
as a special case of the chemical, etc.\ --- every
phenomenon in the last instance as a special case of the
mathematical. But how then can he make the transition
from the general to the particular? Right at the transition
from mathematics to physics the question of
matter comes up, a question that is wholly evaded. In
the same way the particular is constantly presupposed, whereas
according to his doctrine it should always be shown to be composed
% TR: p. 204 the print's „Sammensat“ (capital S in mid-sentence) is a printer's error for „sammensat“; rendered "composed"
of the abstract elements. Because the laws of mathematics are valid
in the whole of existence, all phenomena are not therefore
mere mathematical examples; because the laws of chemistry are not annulled,
but on the contrary confirmed, by the investigation of
organic beings, these do not become mere chemical products.
The fundamental error here is always the confusion of cause
and elementary condition.

By its original principle positivism is essentially
empiricism. The positive was indeed determined in such a way that it coincided
with the factually given, the phenomenal, the relative.
Auguste Comte therefore declares materialism
unscientific, because it sets up an absolute as the final ground
of everything and thereby goes beyond the given. By the whole
opposition in which he places positivism, as that which merely
% --- p. 205 ---
\opage{205}aims at establishing the empirical connection of phenomena,
over against metaphysics, as that which seeks the inner grounds,
the essence of things in and for itself, Comte approaches the skeptical
empiricism that Hume founded. It
was precisely from the denial of the validity of the concept of cause
that Hume set out, and cognition became for him just such a
connection of facts as Comte designates by positive philosophy.
The well-known renewer of Hume's doctrine in our time, John
Stuart Mill, was a disciple of Comte, whose view he
interpreted and developed in an empirical-skeptical direction. But Mill
counts it among Comte's faults that he never admits
open questions. In particular he holds that, because
positive knowledge can only show that in the universe, or rather in the
part of the universe that we know, the phenomena are linked
to one another in a determinate way, it therefore cannot deny
the possibility that this whole lawful existence has its
final ground in an eternal thought, provided only that one does not let
this act in an arbitrary and lawless manner. What opinion
one forms about this depends, according to Mill, on the significance
one ascribes partly to the analogies that can be found in
nature with conscious striving toward ends, partly to
the general traditions of the human race\footnote{\textit{Auguste Comte and Positivism, p. 15.}}. Here
a skepticism is introduced which turns against the very principle of
positivism and reduces it to having only relative and conditioned
validity. But when positivism can no longer maintain
its own absolute validity, then the barrier that separated
it from skepticism has also fallen.

Mill finds the explanation of certain extravagant doctrines
in Comte's later writings in a tendency toward unity and
systematizing which, in his opinion, is peculiar to
French thinkers. This tendency makes itself felt in
Comte's whole fundamental view; it was this that led him
dogmatically to set up the proposition of the explanation of the concrete
% --- p. 206 ---
\opage{206}by the elementary, whereby he passed over into materialism.
This side of his doctrine was most emphasized by his
French followers, especially by the best known among them,
the learned naturalist and philologist Littré. Stuart Mill,
on the other hand, parts from him also in the way in which
the aim and essence of ``positive'' knowledge are more closely conceived.
Whereas Comte seeks to abolish all particular phenomena by
deriving them from their elements, Mill insists that the
ultimate laws of nature that can be reached cannot be
fewer in number than the qualitatively different sensations and
states of consciousness \textit{(sensations or other states of consciousness}»);
the limit must be set where the quantitative and
gradual ends\footnote{\textit{System of logic. 5 ed. II. p. 5.}}. This determination of the limit of
positive knowledge, too, points us in a quite different direction
from the materialistic one.

But neither did Auguste Comte himself hold fast to that materialistic
fundamental principle. When in »\textit{Cours de philosophie
positive}« he came to treat of organic beings, he set up
a determination of the essence of the organism which
recalls the Kantian one. With inorganic beings, he says,
the path of cognition goes from the parts to the whole;
but with organic beings the true cognition
of the parts must be derived from the cognition of the whole.
What is peculiar to the organic being is namely that
it is a closed unity that develops; the explanation of
it must then show the ground of this unity and development.
Analysis, which dissolves the whole into its elementary parts, is
not sufficient here; analytic activity is only the preparation
for the concluding cognition, which rests on a
synthesis that grasps the whole. This now leads Comte to
set up a proposition wholly opposed to that earlier one.
The higher, he now says, explains the lower, not
conversely. To derive and explain the higher from the lower is
% --- p. 207 ---
\opage{207}materialism. This he therefore now rejects as unscientific,
no longer because it goes beyond the empirical,
but because it proceeds solely from the elementary.

The new law Comte sets up here he does not, however,
apply at a point where its application seems precisely to
lie close at hand, the relation between mind and body. Stuart
Mill rightly reproaches him for having no place
for a psychological science, but going straight from physiology
over to sociology (the doctrine of human society).
The reason for this is that he does not acknowledge
a peculiar mental principle, but regards the
mental merely as a product of the bodily. Here he remains
true to his earlier fundamental principle: to derive the
higher from the lower. Spiritual phenomena are at
once more composite and more special than the physiological
and physical ones. When, despite this intricate composition,
a mental unity, an ego, shows itself, this is only an abstract
concept which forms gradually through the feeling of
the harmony and equilibrium which, in the normal state, forms
among the manifold, mutually independent forces
and functions in human nature. »\textit{La notion du
moi, c'est à dire, du consensus universel de l'ensemble de
l'organisme}». It is therefore futile to try to proceed
from the ego, since this expresses only »\textit{un état purement
fictif}»; the true, which lies at the basis, is the elementary
manifoldness.

Here too, however, the point is to hold fast to the reality of the unity
as well as to that of the conditions. What Comte designates
as »\textit{purement fictif}» is the fundamental peculiarity of the human life
of the spirit itself; it does not let itself be dismissed with a
peremptory decree. It is not only in psychological respects that
the concept of the ego, the unity of consciousness, is of decisive importance.
The laws of human cognition point back to
the inner organization of consciousness. When positivism regards
cognition as a mere collecting and summarizing
of factual phenomena and their mutual relations, then
% --- p. 208 ---
\opage{208}this is a purely mechanical and outward-directed activity, which with
necessity points back to an inner ground by which the
mechanical is conditioned and from which the activity proceeds.
And this inner ground can be nothing other than the
collecting and summarizing thought itself, whose laws become
manifest in this collecting and summarizing working.
That the law of thought is in itself one with the law of actuality is
the final presupposition from which all scientific inquiry
must proceed. Only then does cognition gain inner necessity
and true unity. The unscientific character of positivism --- we may here
add: and of materialism --- then consists, in the
last instance, in its not investigating the first principles of science itself,
not setting up a theory of cognition. But neither
\emph{can} it enter into these investigations without
giving itself up; its essence is immediate dogmatism.

Philosophical criticism has to guard the universal
laws and principles of science, especially against attempts
to set up immediately as a scientific total view what
has validity only in a single sphere, from the standpoint of a single
science. Philosophy itself has at times
believed that by virtue of its universal character it could master
the separate sciences; in our day the opposite tendency
prevails. Philosophical criticism turns against both
one-sidednesses. It goes about, as Socrates once did,
and asks those ``who, because they can practice their own art
in a fine manner, also think themselves wisest in other
respects'' how then they can give satisfaction with regard
to the final grounds and concluding results of knowledge and science.
Before this tribunal positivism and
materialism prove their unreliability, --- and it is before this
tribunal that we also summon the author of the interesting
% TR: p. 208 „de interessante Foredrag“ has the plural article, so „Foredrag“ is plural here: "lectures" (the subtitle's singular „Foredrag“ is kept as "Lecture" per the brief)
lectures ``on the relation between mind and body''; he too
``practices his own art so finely'' that he has supposed he could
give satisfaction beyond its limits as well.

\noindent\hspace*{1.5em}(To be continued.)\par
\begin{flushright}\textbf{Harald Høffding.}\end{flushright}
\begin{center}\rule{2.2cm}{0.4pt}\\[-10pt]\rule{2.2cm}{0.4pt}\end{center}

\clearpage
% --- p. 309 ---
\opage{309}\begin{center}\textbf{\large Positivism and Materialism.}\end{center}

\begin{center}\rule{1.8cm}{0.4pt}\end{center}

\begin{flushright}\begin{minipage}{0.6\textwidth}\small\textit{Imperfect as is the science of mind, I do not scruple to affirm, that it is in a considerably more advanced state than the portion of physiology which corresponds to it.}\\ \hfill\textit{John Stuart Mill.}\end{minipage}\end{flushright}

\begin{center}\rule{1.8cm}{0.4pt}\end{center}

\secnum{3.}

It is not accidental that most of the spokesmen of materialism
in our time are physicians and physiologists. The study of the
human organism leads one to see how decisive a
significance the bodily organs have for the whole spiritual
life. La Mettrie became a materialist when, during a violent
fever, he had had occasion to observe in himself the effect of the flow
of blood on the brain; from that moment he regarded
thinking as a mere consequence of the organization of ``the bodily machine''.
According to \textit{Système de la nature}, medicine
solves the riddle of the human heart. Medicine has therefore
been called the theology of materialism.

Yet it is far from the case that physiological study
must of necessity lead to materialistic consequences.
One need only look to men such as Flourens and
Claude Bernard, Johannes Müller and Virchow to understand
this. But in general those on the physiological
% --- p. 310 ---
\opage{310}side will speak out in a positivist direction. That Maudsley in particular
stands on the standpoint of positivism can easily be shown.
True, he expressly invokes ``the spirit of positivism''
in only a single place; but that his whole way of looking at things is
governed by this spirit is seen both from what he combats,
from what he declares inaccessible to cognition, and
from what he takes as his point of departure.

Science, he says, has inherited from theology the
habit of regarding the mind as an incorporeal being that
could not be apprehended through the senses. Thence arose ``the
metaphysical philosophy of mind'', which did not take due account of
the physical conditions under which mental life develops.
One began where one ought to end, setting out
the philosophy of mind according to the dubious revelations of self-consciousness,
whereas it was precisely these revelations that
were to be explained. The body was underestimated and despised,
regarded as something one had to be ashamed of. Maudsley's
efforts are therefore directed at freeing the problems
from the fog in which metaphysics is supposed to have shrouded them.

In opposition to this it is insisted that the question of
the nature of the mind is such that science cannot
take hold of it at all, but that we are, and in all likelihood will always
remain, completely ignorant in this respect. We
can only feebly grope after the laws of its activity.

In what way, then, can its laws be found? Once
the mind is no longer regarded as a non-sensible being,
nothing stands in the way of applying the inductive method
in the investigation of the mind just as in that of
all other natural phenomena. We begin, then, with the simplest
facts, namely the bodily conditions and expressions
of the mind's activity. Anyone who would enrich our knowledge
of mental life must in our day first and foremost explore
the bodily conditions of its appearance in the healthy
and in the diseased state. What we then gain by considering
the simpler facts we use to explain the
more intricate and difficult ones. The higher functions of consciousness
% --- p. 311 ---
\opage{311}are developed from the lower, just as a peculiar and intricate
structure is a development of a simpler and more general one.
From the concrete we derive the abstract, from the uncompounded
the compounded. The life of consciousness is the sum or product of
a stream of forces to which every bodily organ
sends its contributions, whether these are noticed or not.
The sum must be resolved into its addends, the product into its factors,
in order to be understood.

One sees that Maudsley shares Comte's conception of
the history and character of science. The theological, the
metaphysical, the positive period --- this tripartition stands out
clearly in what has just been cited. And just as Comte declared
the essence of things in and for itself, or, as he mostly
puts it, their causes, to be absolutely uncognizable, so
Maudsley places the question of the nature of the mind
outside science. What then remains as
the task of science is to investigate the conditions, the laws
of the mutual relations of things. Finally, Maudsley also
agrees with Comte in the general proposition that the higher
is to be derived from the lower, the elementary. But the manner
in which Maudsley has conceived the principles of positivism, and
the classic exposition of examples with which he illustrates
them, give his lecture a peculiar interest. --- We shall dwell
first on the general principles before we pass on to
the particulars.

The polemic against ``metaphysics'' is kept, as in Comte,
rather indefinite, and it is therefore not easy to go more closely
into it. On the one hand it seems as though
metaphysics were meant to designate any answer whatsoever to
the question of the essence and nature of things, thus here
of the nature of the mind. But on the other hand
Maudsley designates by that name in particular the spiritualist conception,
% TR: pp. 311--312 render Sjæl as "soul" (4x): the spiritualist doctrine of an immaterial soul, Plato's psyche and Lemoine's thesis that the soul itself cannot be ill; elsewhere Sjæl = mind (independent review)
which regards the soul as an immaterial being that is in
itself independent of the bodily conditions. But
even then it is doubtful whether his polemic really
hits the mark. It is precisely the spiritualists who have been most atten%
% --- p. 312 ---
\opage{312}tive to the significance of the bodily conditions; the higher
they raised the essence of the soul above the physical world, the
stronger the influence that the bodily conditions
exerted on it had to become. On this point spiritualists and
materialists could splendidly agree. Plato, that father of
spiritualism, teaches expressly that it is the constitution of the body
that produces the diseases of the soul, that ``no harmony and
disharmony has greater influence on health and disease,
virtue and vice, than that between the soul and the body''; --- indeed, he
% TR: p. 312 the print has no comma after „Materialister“ at the line end (logged in the transcription); the English needs none here, and the sense is unaffected
goes as far as any of the most extreme materialists
when he says that no one is bad by his own will, but that the ``bad
become bad through a diseased bodily constitution and through a bad upbringing''\footnote{Timæus. (Heise's trans., pp. 138--141).}. And as regards the particular question of insanity,
which is the one to which Maudsley ties his investigation
of the relation between mind and body, it is precisely
spiritualist thinkers (such as, e.g.,
the Frenchman Albert Lemoine) who in the most recent times have maintained that it must necessarily
have a bodily cause, because they deny that the soul in
itself can be diseased. --- If, on the other hand, the polemic against metaphysics
is to apply to every attempt at a solution of the problem
of the nature of the mind, it applies to materialism just as much
as to spiritualism and establishes an empirical
skepticism. We have seen above that this skeptical element
was characteristic of the point of departure of positivism.

But Maudsley holds fast to this skepticism as little as Auguste Comte
did. Through the manner in which he applies the
inductive method he arrives at an actual answer to
the problem he at first rejected, ``really takes hold of'' the question
of the nature of the mind --- and to that extent becomes
a metaphysician himself. There is no doubt that he is perfectly
right to apply the inductive method in the investigation
of mental life; but it depends on how the presuppositions
for the application of the method are originally set up: an error
here must necessarily, despite the most acute and exact
% --- p. 313 ---
\opage{313}carrying through of the method, produce an error in the result.
Now Maudsley commits from the beginning the error of regarding
the elementary conditions of mental life and its
rudimentary forms as the concrete and simple, from which
consciousness in its various forms emerges by composition;
the higher mental life is then designated, in relation to the
elementary, as an abstraction. And yet it becomes clear from
Maudsley's own exposition that it is the reverse
relation that is the true one. ``It becomes necessary,'' he says,
``to separate out the lower and more elementary of these
functions and from the concrete to derive, as well as possible,
how the abstraction is to be understood.'' But what
is this separating out other than an abstracting? The single
function is drawn out of the whole connection in which
it has its peculiar existence. The elementary that is
separated out in this way becomes abstract. What Maudsley
calls ``the abstraction'' is the total connection in which
all single functions work, and which embraces and guides them,
--- a connection which at its highest stage of development is
consciousness itself, the mind, insofar as it, to use
Maudsley's own expression, ``gathers and comprehends the bodily
life.'' This gathering and comprehending presupposes a
peculiar activity and proves that the mind is not merely an
abstraction or a sum of forces. Through the total
connection, which conditions the elementary functions just as much
as these condition it, mental life proves itself to be
something concrete and real in itself. And in thinking the active
unity of the mind comes forth still more clearly. Every act of consciousness is
in itself a single, indivisible act. --- By turning attention away
from this real unity of the mind and of consciousness,
Maudsley has made the task unduly easy for himself,
just as easy as he began by making it difficult when
he declared the nature of the mind to be altogether incomprehensible. And
not only this. Through the turn he has given
the investigation, he has passed over into the materialistic conception
of mental life. For it is the peculiarity of this conception to
% --- p. 314 ---
\opage{314}posit the mental as the merely derived (sum, product,
abstraction) and conditioned in relation to the bodily, and
not to grant the former a peculiar essence of its own.

It is, as has been said, not the inductive method itself that
leads to this result, but the presuppositions that are laid at
the basis in its application, and into which
determinations of deeply far-reaching significance are smuggled. Maudsley
draws attention to ``the realm of little things'', the
infinitely fine nuances and movements that easily escape the investigator's
eye and yet can be of great significance. But
in a logical respect he has himself overlooked the ``little things''
which, hidden in the presuppositions, appear in the consequences
as very big things. At the outset he commits a
logical, or if one will a metaphysical, mistake, which exerts
its influence on the whole interpretation of the physiological
results, but for which these are not to blame. The physiological
investigations do not lose their value because the presuppositions
for the concluding cognition, --- a cognition that goes
beyond the limit of physiology, --- are false. --- We shall
now pass from the more general considerations to the particular,
in order here too to investigate the relation between
positive inquiry and the results at which it believes it
has arrived.

It is from Pflüger's and Flourens's well-known experiments that
Maudsley takes his point of departure, and it is the interpretation
of these experiments that is at issue. A frog that has
lost its head draws away the leg that one pinches, and
hops away if the irritation continues; if one cuts off the foot
of the leg with which it rubs off a corrosive acid,
it takes, after several vain attempts, the other leg
to its aid. A pigeon from which one removes the hemispheres
% TR: p. 314 „den store Hjernes Halvkugler“ (the hemispheres of the great brain) rendered "the hemispheres of the cerebrum", here and on p. 315
of the cerebrum loses, so to speak, its initiative, but nevertheless flies
when one throws it into the air, closes its eyes against the light,
and swallows the food one pushes into its mouth.
Now are these so-called reflex movements or automatic
% TR: p. 314 the print's „ellér“ (a foul-case é, logged in the transcription) is „eller“; rendered "or"
actions mental or merely bodily? To ascribe them to the mind
% --- p. 315 ---
\opage{315}would in Maudsley's opinion be a kind of anthropomorphism:
it would ``smack of the old, bad tendency
that has done so much harm in philosophy, and that
consists in explaining the facts of nature by
what we feel in our own mind.'' On the other hand he infers from
those experiments that actions with a definite aim can be
altogether unconscious: ``The animal has had deposited in its spinal cord
the capacity to perform such movements for self-preservation
as are given to it as a part of its nature, and without
which it could scarcely live a day''. In man, too,
something similar shows itself. Not only are sneezing and coughing
a reflex movement, a movement that, independently of the will,
serves to remove what is harmful to the organism; but
even the acts of will proper presuppose for their execution
automatic actions proceeding from the motor centers of the spinal cord;
for the how of the movement does not, after all, fall under the
conscious will. And when a person repeatedly does
something, the spinal cord is, as it were, trained to
act on its own, so that the action is performed quite
instinctively and automatically while consciousness itself is busy
with entirely different things; so it is with many daily actions
(dressing and undressing and the like). According to
Maudsley it is no small part that is subtracted, by the doctrine of reflex movements,
from the phenomena we are accustomed to regard
as mental.

But the conclusion here does not seem to be as exact
as the experiments from which it is drawn. So long as
the one possibility stands beside the other, one
cannot, after all, say that the answer is positive. Cuvier and Johannes
Müller drew with full right the opposite conclusion
from Flourens's experiments, namely that the animal, after having lost
the hemispheres of the cerebrum, did indeed become dull and torpid, but
nevertheless showed clear signs of actual sensation. How
will Maudsley be able to refute this interpretation (which,
remarkably enough, he does not even mention), since it has the signs
on its side? All the more so since he would not be able to avoid
% --- p. 316 ---
\opage{316}coming to something similar if one were to demand of him
a closer explanation of the mystical expressions he would put
in place of sensation and will: ``capacity for self-preservation,
to feel \textit{(sic!)} and avoid what is harmful, as
also to feel and strive \textit{(sic)} after what is beneficial.''

The fact is that Maudsley, through the ``separating out'' discussed above,
which he did not see was an abstraction, has set a
sharp distinction between the mental-conscious and the organic,
unconscious life. But such a distinction cannot be maintained
without manifold points of transition. Within sensation
itself there is an infinite scale of stages leading from
unconscious to conscious life. Conscious life therefore cannot be derived from
unconscious life except by being presupposed in it as disposition and real possibility.
There is an inner relation between the mental principle
and its conditions, such that from the very first these enclose
the principle within themselves, in order step by step to let it
develop as they develop themselves. The degree of consciousness
can thus vary; but when, as Maudsley himself concedes,
one has before one actions ``which in actuality
bear the stamp of preceding consciousness and will'', there is
no reason to think that consciousness should be entirely
lacking\footnote{\emph{Virchow} uses the concept ``unconscious sensation'' to explain reflex movements. (Ueber das Rückenmark. Berlin 1871 \textit{p.} 25). This indicates the right intermediate determination, when one remembers that sensation can never stand in complete opposition to consciousness. --- He points out, moreover, that the influences that call forth reflex movements are propagated through the sensory nerves, and that there are, besides, also reflex movements caused by \emph{conscious} sensations (e.g., contraction of the eyes on account of too dazzling a light).}. Moreover, the automatic actions in man do, after all, exert
a considerable influence on the life of consciousness itself;
but how could that happen if they took place entirely outside
consciousness? The mind cannot ``gather and comprehend''
what falls outside its domain.

% --- p. 317 ---
\opage{317}Maudsley does not stop at the mere reflex movements.
What he believes he has gained here as a result,
he seeks to apply to the higher mental functions.
Despite the difference of the functions, he nevertheless holds that the properties
and laws are essentially the same as in the lower
nerve centers: ``Impressions are received, upon these impressions there follows
a reaction, and the effects both of the former and of
the latter are ordered and preserved.'' The impressions are the physiological
conditions of \emph{representations}. ``The sensation of representations
is \emph{feeling}; for by feeling I understand
the peculiar receptivity of the nervous substance to representations.
That these are ordered and preserved is \emph{memory}, and in the
reaction they call forth, the \emph{will} is given.''

Here, then, we have in brief what is to take the
place of the ``metaphysical philosophy of mind'' that is supposed to
have done so much harm. This much is certain, that these
determinations are not metaphysical; it is another matter whether
they therefore have less of the nebulous that, in Maudsley's
opinion, is peculiar to metaphysics.

The first thing that strikes the eye in those determinations
is that it is not said at all what a representation is. It is
only stated that the impressions are its physiological condition,
and with that its essence is regarded as explained. This
superficiality has its ground in the fact that, in order to explain
the essence of representation, one must necessarily presuppose consciousness;
but this Maudsley, who wants precisely to prove that consciousness
arises as a sum or composition, cannot do.
He therefore keeps to the empirical conditions and regards
them as a sufficient explanation\footnote{\textit{``It is one question what a thing is, and another what it depends on''.} \emph{\textit{Stuart Mill}}. \textit{System of logic. II p. 7 (5th. ed.)}}. ---The matter does not become
clearer because it is said in another place that ``representations
are acquired through the senses, mixed, bound together, and grouped
---and these organic combinations are the physiological
% --- p. 318 ---
\opage{318}conditions of the highest forms of our spiritual activity:
reflecting, reasoning, and judging.'' Since we
know beforehand only so much about representations, that they depend
on physiological conditions, it is not easy to say what a
mixing, binding together, and grouping of them is supposed to mean.
Yet---when we further hear that these are ``organic combinations''
which in turn become ``physiological conditions,'' then
the representations themselves can be nothing other than organic
processes, and thus do not merely \emph{depend} on such processes, since such a
dependence does not exclude essential difference. With this the determination
of feeling would agree: ``the peculiar
receptivity of the nervous substance to representations''---the representations must
be something bodily in order to be able to affect the nervous substance.
But the same difficulty recurs here: surely no one can
seriously hold that from the receptivity of the nervous substance
there follows the least thing toward the explanation of a phenomenon
such as feeling, which is an inner state and presupposes a
relation of the being to itself.

In the face of these difficulties Maudsley has no right
to appeal to the uncognizability of the essence of the mind.
% TR: p. 318 printer's error: the last letter of „tidligere“ is a wrong sort (Danish as printed, logged there); rendered in its evident sense, "earlier"
For we have already seen earlier that he himself does not hold fast to
this uncognizability. Nor does he do so in his
physiological determinations. When it is said, e.g.: feeling
is the receptivity of the nervous substance, etc., then this
proposition appears with the same positive authority as his
other doctrines. On the contrary, what follows from those difficulties is
the necessity, in investigating mental phenomena,
of taking one's point of departure not merely in the physiological
conditions but within the mental itself, in consciousness.
And this is not a new and foreign point of departure. The life of consciousness,
with its peculiar laws, points back, both through
its successive development, its crises and diseases, and through
its very end, the appropriation and development of the content
of existence, to real, physical conditions as its
foundation. The psychological in itself points toward the
physiological. If, on the other hand, the physiologist does not have a clear conscious%
% --- p. 319 ---
\opage{319}ness of the peculiarity of the psychological, the kernel of the question
will escape him.

That this consciousness is lacking in Maudsley is seen not
least from his doctrine of memory. According to him, memory consists
in an ordering and preserving of representations. But
it is, he continues, a metaphysical error that memory
should be peculiar to the mind. The acquired functions of the spinal cord
and of the sensory ganglia cannot
go on without memory; these organs would become
idiotic if there were no memory in them. ``In every
nerve cell there is memory, and, what is more, it is found
in every organic element of the body. The poison of smallpox or syphilis
sets its mark on the body for the rest
of life; we may forget it, but it does not forget us.
The way in which the scar from a cut on the finger
of a child persists and grows with the child's body
proves that the organic constituents of the finger remember the
change they have undergone. Memory is an ordered
preservation of the effects that the impressions have called forth.''
But where are we now headed? Imperceptibly the first, psychological
determination of memory is lost from sight, and all
ordering and preserving of effects of impressions is made into
memory. No account is taken of the essential difference
between the way in which representations are stored
in consciousness and the way in which effects are preserved in the organism.
A name that is carved in the bark of a tree
is surely not preserved by the tree in the way it is preserved by the one
whose ``dear, constant thought'' it is. One of two things: either
Maudsley must here, in the manner of a poet, make the tree into an ensouled
% TR: p. 319 „besjælet“ rendered "ensouled" (Sjæl = mind elsewhere; "minded" will not carry the poetic sense)
being with remembrance and sympathy, or else he must
make consciousness into a purely mechanical container. What
produces the unclarity is that memory is unthinkable
unless the organic elements retain the movement corresponding to the remembered
representation, or rather its possibility.
But this organic movement is not the representation itself,
only its condition. In the organic element the
% --- p. 320 ---
\opage{320}movement corresponding to the representation is ``remembered'' in the same way as
the scar in the boy's finger, the name in the tree's bark; in consciousness
the representation is remembered as a moment in the whole
world of thought that is the content of consciousness, and it is remembered
with an intensity that depends on its significance in this
content. Such words as impression, order, collecting
therefore acquire a different meaning according to the side from which
one views the matter. But the positivist, who overlooks this,
comes not only to speak physiologically where he ought to speak
psychologically, but also psychologically where he ought to speak physiologically
---and thereby goes back from the ``positive'' standpoint
to the ``metaphysical,'' indeed perhaps even to the ``theological,''
in that he explains natural phenomena after the analogy of his own
essence.

No more than memory is the will, according to Maudsley,
a peculiarly mental phenomenon. For it is, says
he, a fundamental peculiarity of all organized tissues
that they strive after what is beneficial and shun what
is harmful to them. The impulses that arise from
this property of the tissues we call manifestations of will when
they are guided by reason. ---Apart from the difficulty, already
mentioned earlier, of conceiving what it is supposed to mean that
an organized tissue strives, shuns, and impels---there is
this peculiarity about this determination of the will, that the
decisive factor is introduced as a condition that is not explained.
What distinguishes the will from mere drive is, to be
sure, that it is guided by reason; but then it must first
and foremost be shown of what kind this relation between reason
and the drive-impulse is, and how it is possible.
How can the impulses at once be called forth by
the properties of the organized tissues and ``according to reflection
be guided by reason''? This question, which psychologically
is the main question, Maudsley does not raise at all, any more
than we get to know what he understands by reason. He thereby
comes to lack a real standard for
% TR: p. 320/321 the page turns inside „Villies-|yttringer“ (manifestations of will); the English turns before "will"
the relation between the reflex movements and the proper manifestations of%
% --- p. 321 ---
\opage{321} will; for when, in the case of the latter, one sets aside the element of reason
or consciousness, they certainly do not become
different from the former.

One can, it is further said, speak physiologically only
of manifestations of will; the \emph{will}, as something different
from the individual actions, is a metaphysical abstraction.
---Whether there really is anyone who has wanted to place
the will outside the actions must surely be doubted; Maudsley's
polemic is here, as on other points, not free from itself
producing the opponent it combats. The same polemic
could in any case be directed against him
himself. The fundamental peculiarity he assumes in the
organized tissues likewise becomes, when his manner of consideration is applied,
merely a ``metaphysical abstraction''; for it manifests
itself only under the influence of determinate agencies. But
what does one understand by will other than the fundamental peculiarity
of mental life to which the manifestations of will point
back? Positive inquiry, which keeps to the factual
effects in order through them to come to know the essence of things,
cannot, after all, as Maudsley's own example shows, avoid
leading the effects back to the essence as possibilities
contained in it. This is perhaps metaphysical; but then it
is a metaphysics to which physics itself leads us. If
one wants to say, e.g., that we do not know the essence of electricity,
this must be restricted to this: that we know electricity
insofar as we know how it acts under determinate circumstances.
We know, after all, that the electric current forms
a link between the different natural activities, since
by it one can produce light and heat, magnetic and
chemical activities, and conversely the current from these; and one
can measure the electric current by its chemical or magnetic
effect\footnote{C. V. \emph{Holten}: Om Virksomhedernes Omdannelse i Naturen. 1868. (University program). p. 19.}. Should this not be a contribution to the
cognition of the nature of electricity? If one seeks the essence
% --- p. 322 ---
\opage{322}or nature outside the manifestations, one is guilty
of a false metaphysics. We have indeed witnessed the strange spectacle that a
thinker like Stuart Mill has written a book on the position of woman
in society that contains the most ingenious and most
apt characterization of woman---but in which he nevertheless
maintains from first to last that we know nothing of the nature
of woman!

\begin{center}\rule{1.8cm}{0.4pt}\end{center}

\secnum{4.}

Partly as a means of further elucidating what he
has developed about the relation between mind and body in general,
partly as a particular application of it, Maudsley
enters upon an investigation of the essence of mental illnesses
and their relation to other nervous disorders. These
investigations are the most interesting and most instructive in
his lecture; they are also those in which he most avoids
an unfortunate overstepping of the boundary of his science and thereby also
of his competence. When we nevertheless attempt here too
to apply a critique, it must, even more than
on other points, be kept in mind that this critique proceeds solely from a
psychological point of view and therefore cannot detract in any way
from the significance that Maudsley's investigations may have in
a medical respect.

The usual conception and classification of mental illnesses
does not satisfy Maudsley, because it proceeds exclusively
from a psychological standpoint. The true scientific
spirit, ``the spirit of positivity,'' is in
his opinion satisfied only when one holds fast to the fact that these
disorders rest on the same conditions as other diseases.
Since mental life is so closely bound to and dependent
on the bodily functions, mental illness cannot
be a purely mental phenomenon either; it is not a
disease of a metaphysical concept, but a nervous disease
with predominantly mental symptoms. This view is
of course not new; among us, for instance, F. G.
% --- p. 323 ---
% TR: p. 323 Howitz's title kept in Danish as printed (``on insanity and accountability'')
\opage{323}Howitz already urged it in his treatise ``om Afsindighed og
Tilregnelse'' (1824), where insanity is defined as ``a
restriction of reason or of the use of reason owing to
a disease in the material organs of the mind's activity,''
as ``brain sickness,'' and where weight is therefore laid
on the double (mental and bodily) investigation and
determination of it. But what is peculiar to Maudsley's
investigation is the closer grounding and the further
consequences.

One does not know at all the movements in the parts of the brain that
are the physical conditions of mental activity, or
the chemical changes that accompany them, any more than
one can demonstrate any difference between the nerve elements of the tired, exhausted
brain and of the revived, refreshed one.
The means of investigation do not suffice in the face of changes that
are finer than the microscopic. ``We might almost as
well set about studying the anatomy of a gnat with a
telescope.'' Thus the brains of deceased mentally ill persons
have been found entirely normal. But according to Maudsley one must
not conclude from this that the change is not there; it may well
exist, even if it belongs to ``the realm of infinitely small
things,'' which is inaccessible to our senses. By virtue of
the close connection between mind and body and of the
bodily symptoms that accompany mental illness, a disease in the brain
\emph{must} be assumed.

The justification of this inference, however, does not stand quite
unchallenged. One must, in the first place, wonder that
Maudsley takes no account of the cases, which have been
established, in which both anterior lobes of the cerebrum
were diseased and destroyed to a considerable degree, and in which nevertheless
% TR: p. 323 footnote: Lange's title kept as printed, „Geschichte der Materialismus“ (der for des)
not the slightest disturbance of the intelligence was observed\footnote{\emph{Longet} cites two such cases. See \emph{Lange}: Geschichte der Materialismus. p. 431.}.
And next, one will be able to set against Maudsley's line of thought
another, apparently more scientifically sober one, that of
% --- p. 324 ---
\opage{324}the French physiologist Leuret. Certainly, says the latter,
I have no right to conclude that there is no change
because I cannot see any; but with the same caution
I must guard against concluding that there \emph{is} one. When
the brain of an insane person appears to be healthy, I do not conclude
that this brain is diseased, but I remain in doubt until
the truth is proven. And if the cases in which the brain
appears to be healthy are precisely those in which the insanity is not
accompanied by bodily symptoms but is an aberration of
the understanding and the passions, while the cases in which
the brain is altered are those in which paralysis,
dullness, and sleeplessness have occurred, then I ascribe the brain damage
to these circumstances, and the cause of the derangement of mind
remains unknown\footnote{\textit{Revue des deux mondes. 15 juillet 1865. p. 417.}}.

This much at least seems to follow from this, that
drawing too hasty conclusions in this field is attended with the greatest
hazards. But in any case
---the correctness of the claim that there is always a bodily cause granted---
has one thereby gained a closer insight into
the essence of mental illness, and must this not still
be understood primarily from the psychological side? Maudsley himself
remarks that the mental symptoms are the predominant ones.
But how can one demonstrate an inner connection between
them and the bodily symptoms? It has rightly been
said that when a clock strikes wrongly because a wheel is
broken, it does not follow that it was this wheel that
struck;---we do not know the nature or operation of a locomotive
because we know that it cannot do service when it is
broken. The connection between the mental and bodily
symptoms becomes a purely empirical one, whereas within
mental life we can follow insanity step by step
from the first faint signs to the demonically overwhelming
power with which it appears at its height. Here
% --- p. 325 ---
\opage{325}the laws of the association of ideas, of thought, and of character come
to our aid --- but the study of these laws is a psychological study.
The physician of the mentally ill must therefore above all be a practical
psychologist.

Maudsley does not seek the bodily cause of insanity in
an influence from without, through a spiritual shock or bodily
% TR: p. 325 „aandelig Rystelse“ rendered "spiritual shock" per the glossary (Aand = spirit); the sense is a mental or emotional shock (also l. below, „samme aandelige Rystelse“)
injury; such an influence is only the occasion.
The real root of mental illness lies deeper, in an innate
disposition, in the organization of the individual's nervous system. When
two people are exposed to the same spiritual shock, and the
one of them becomes insane while the other only for a
moment loses his equilibrium, the reason for this can only
be that in the former there is a factor in play which is not
found in the latter --- some innate nervous disease or other.
Should one hold that this need not after all be innate,
but could have formed under the influence of the circumstances of life,
Maudsley rejects this in the following way:
``Of all the erroneous conceptions of the mind which metaphysics
has begotten or been complicit in, there is perhaps none that
is more false than that which tacitly assumes or expressly
declares that all men are born with equally good conditions
for the development of a mental life, while circumstances
and education condition the differences in later
development''. Poor metaphysics, which now, besides its
own sins, must also take the offenses of materialism upon
its shoulders, and is thus here accused of the very opposite
of what was brought forward against it above! Still, it is not
ours to defend it here. We shall instead see what
Maudsley sets up in opposition to what he lays
to its charge. ``Education can certainly accomplish much, and
the circumstances of life can also accomplish much; but
we must not forget that the foundation on which the yield of
education must rest is not acquired but inherited. No one
can escape being ruled by his organism, no one can
flee the fate he is born with''. The innate is
thus inherited, and the inherited fund of forces that man
% --- p. 326 ---
\opage{326}brings with him to the battle of life conditions the success with which
he fights it out. The life of the mind, like all life, is placed
among forces that always seek to annihilate it, and against
which it must therefore constantly be on a war footing. He who
can master circumstances gets ahead; the weak man
is mastered by them and collapses --- and one of the ways
in which his misfortunes can show themselves is precisely as mental illness.
``Thus it comes about that mental trials,
which serve to strengthen a vigorous nature, break down a weak nature
that cannot offer due resistance; the circumstance that
a moral cause can bring on a mental illness betrays a
conspiracy from within with the unfavorable outward circumstances.''

We inherit as a natural gift what our forefathers
have acquired, both the good and the evil. The moral
sense arises in that what the one generation attains
of conscious cognition of what is beneficial or harmful
to humanity is, by virtue of the organization of the nervous system (cf.
above on memory), passed down by inheritance to the next generation
as ``a natural or instinctive susceptibility that belongs to
his brain when it has its proper structure''. What has been
acquired piece by piece is put together so closely that it finally
seems to constitute one single, indivisible faculty. By following
the history of humanity back, one would finally, at the starting point,
find ``a state which, whether one would
call it apelike or human, lacked any real moral
feeling.'' The physiological proof of this is supposed to lie in the
progressive development of the human brain; the part of the
tissue that has been added represents the experiences
and recollections of the race, which have been incorporated in us. But
it is only when the normal presuppositions are present that
the moral sense is inherited. If the structure of the brain is not
the proper one, if the individual has ``the insane temperament'', then
% TR: p. 326 „det sindssyge Temperament“ rendered "the insane temperament" (Sindssyg = mentally ill in the glossary; "mentally ill temperament" will not do as a term)
it is precisely that sense, the highest flower of culture, that is lost
first. If this abnormality persists through several generations,
idiocy appears. He who is born an idiot has
% --- p. 327 ---
% TR: p. 327 the print's „theoride“ (for therioide/theroide) rendered ``theroid'' (Maudsley's term, 'beast-like')
\opage{327}lost his birthright as a human being and sinks back to
the animal stage. Of particular interest here are the ``theroid
idiots'', in whom not only animal instincts but also an
animal-like appearance occur --- traits which are as it were
echoes from a vanished time, testimony to a kinship that
man has almost outgrown and is too proud to acknowledge.
In such failed and degenerate individuals, then, the result
that was won is lost, and its original, unworked
element comes forth again. Thus the psychology of mental illnesses
confirms Darwin's view, pointing back
to it as to the only thorough explanation.

It is with a certain necessity that Maudsley concludes
his investigations by referring to Darwinism, this
bold, concluding view of modern natural science.
But it is not thereby said that this reference is one with an
explanation. Darwinism has a certain indeterminacy in its
essence; it can be understood from two sides. That all organic
forms have developed from one primordial form can namely either
mean that all those forms are only expressions of the different
ways in which the primordial form was mechanically affected and
modified by the outer conditions, or that they were originally
contained as dispositions and possibilities in the primordial form, so that
this form encompassed in itself, as in a germ, the whole world of
living beings. In the latter case Darwinism becomes no
particularly positivistic or materialistic doctrine. For the higher
forms are then so disposed from the first and, as dispositions, have an
ideal superiority over the conditions, which become only the means
of their actualization. And if one notes that the
real main factor in Darwin is not the outer conditions
but what he calls the struggle for life, it seems
reasonable that this latter conception of Darwinism must
be the right one. For a struggle for life there is required an actual
principle of life, where energy can be concentrated, from which
the struggle can be waged. There must be something for whose preservation
and development the struggle is waged; without such an \emph{original} good
the struggle would have led to nothing. Out of nothing comes nothing.

% --- p. 328 ---
\opage{328}But Maudsley starts precisely from an original nothing,
from a zero point. The development of the race is supposed to have begun at
a stage where any real moral feeling was lacking. Only
gradually are the fragments gathered from which the moral
sense is put together in such a way that it can finally appear
as an indivisible faculty. But the very appropriation and working-up
of these fragments presupposes an activity, a
laboring, a ``gathering and comprehending'' --- in short, a working
whose possibility we cannot see if there is not
something originally at work. In psychological respects
Maudsley starts from the view that the English philosophers
of the present day (Stuart Mill, Herbert Spencer) have developed
as an inheritance from their great predecessor in the last century,
David Hume --- the view, namely, according to which the laws and the
whole development of the life of consciousness rest on the associations of ideas,
according to which the empirically received representations group
themselves. The representations, once grouped and received,
gradually become a lasting part of the mental constitution
and pass by inheritance from the one generation to the next, so that
for following generations they present themselves as truths valid in themselves
or assert themselves as immediate instincts.
But this doctrine of the growth of the understanding (\textit{growth of intelligence})
overlooks that the essential, lawful connection
which is peculiar to the content of thinking consciousness
cannot have arisen through a mere accumulation of representations
or through the constancy of impressions from without. The
inner lawfulness is the fruit of a labor of thought which necessarily
presupposes a laboring thought, whose essence is
one with the universal reason in existence. Without this, the
ultimate principle of all human cognition, all
knowledge would be impossible. And besides, since the immediate truths
in the course of time become objects of criticism, since consciousness
can thus turn against its own work in order to annihilate
it, there must be an inner unity in the self, if it is not
to perish under such criticism but is to be able to
work its way forward to new expressions of the truth.

% --- p. 329 ---
\opage{329}The moral sense is not one with what at a particular
time and place constitutes its content. Maudsley here
himself becomes guilty of what he reproaches ``metaphysics'' with:
ascribing too much to the outward. The properly moral is
essentially a relation between self-consciousness on the one
side and its content and conditions on the other,
but it does not consist in this or that particular content. The Australian negro
therefore has a moral sense just as well as the European most
developed in intelligence. But the positivistic
and materialistic way of viewing becomes guilty of much
superficiality with respect to such questions. When,
for example, one finds among a savage people no belief in
particular gods or any cult properly so called, the conclusion is
at once drawn that this people lacks religiousness, without
one's investigating whether the religious does not express itself
in another way among this people, whether it does not, for example,
merge with the relation of piety toward chieftains and fathers, who thus
become a kind of earthly gods. Or one denies
certain savage peoples the capacity for cognition because one does not find
general, scientific propositions expressed in their language.

% TR: p. 329 the print opens the quotation before „Arvelighed“ with a single low mark (‚) and closes it with a double one; rendered as a normal pair ``heredity''
Further, there hides beneath the expression ``heredity''
an ambiguity. A content of consciousness cannot be directly
inherited; it must be appropriated anew by each generation. Here
there can only be talk of immediate inheritance insofar as the growing
generation, through education, takes up without criticism
what the preceding one hands over to it; even after criticism has
awakened, it can be difficult to free oneself from the immediately
received impressions. What, on the other hand, is directly innate,
as a part of the individual's original constitution,
can only be more or less favorable dispositions and
drives. As the essential thing about the cultivated moral
sense is an ethical How, so the essential thing about
the innate possibility is an organic-psychic How. It is
really only here that Darwin's theory finds its application
--- an application that lies plain to see, and
which Plato already made in his ``Republic'', where he
% --- p. 330 ---
% TR: p. 330 footnote: „Platons Stat“ / „Staten“ (Plato's dialogue under its Danish name) rendered "Republic"; „Heises Overs.“ = Heise's translation
\opage{330}in a sense appears as Darwin's forerunner\footnote{Plato's Republic. Book 5. (Heise's transl. of the Republic. Part 2, pp. 24--27).}. The innate
possibility does not lead directly to the goal; it presupposes
the individual's own striving and is not actualized if the individual
will not take the matter into his own hands. Admittedly the
higher development of energy acts back on the organic-psychic
foundation that passes by inheritance; this is plainly seen in the noble
families. But still it is not the happily developed brain;
it is only the conditions for its happy development that
can be inherited.

Maudsley does not apply the same way of viewing
to the race as to the individual. As regards the latter, he maintains
the significance of the innate foundation in opposition to the
effect of outer influences; but in the case of the whole race he thinks
he does not need such an original foundation; there he ascribes
an overwhelming significance to the outer influences.
And yet one will hardly be able to maintain such an opposed
way of viewing; one must necessarily apply a
kindred standard in both places, as surely as the essence of the race
cannot be heterogeneous with that of the individual. The German
materialist Büchner is more consistent here than the English
positivist; he will not even concede a peculiar, innate
foundation for the life of consciousness, for fear that
``innate ideas'' should be smuggled in under this form. If
Maudsley thinks that on grounds of experience he must hold to the significance
of the innate as regards the individual, he must
also do so as regards the race.

But the innate cannot in any case be said without
further ado to be one with the individual's fate, so that
this fate is settled at birth. Fate in the true sense,
the ethical fate, the result of the interaction between
the individual's conscious striving and the inner and outer conditions,
cannot be innate. Conscious striving is, even if
% --- p. 331 ---
\opage{331}it must succumb to unfavorable conditions as to
the power of natural necessity, an ethical reality. Only the purely
idiotic existence is immediately determined by the innate
foundation; but this existence then also lacks the proper
% TR: p. 331 „Karakteermærke“ (so printed) rendered "distinguishing mark"; no translation issue
distinguishing mark of the human. Maudsley rightly declares
that the idiot is robbed of his birthright as a human being,
since he is born with such a deficiency of brain that he cannot
perform any mental function at all, or can do so only
imperfectly. But how one must be astonished on reading his
following words: ``He (the idiot) himself bears no blame for
this visitation, since he has committed no other
offense than that of having his share in human
sinfulness.'' What can this mean, if it is not
a mere remnant of a theological metaphysics? To have a share
in sinfulness and yet not to have committed any offense,
indeed even to lack the possibility of it, since one is from the very outset
robbed of one's birthright as a human being! Sinfulness can
no more be inherited than the moral sense can. Such
a way of viewing leads to a complete confusion
of medical and moral concepts. The counterpart
to the idiot's sinfulness is formed by the conception of crime
as mental illness or epilepsy.

And yet the latter conception is more correct than the former;
it requires only a closer determination. One must not here
confuse the legal and the moral standpoint.
The crime can well be present legally without being
so morally. One can imagine an existence which, even if it
has not lost its birthright as a human being, has nevertheless
arisen and developed under circumstances that, with the power of a natural
necessity, must have forced back everything higher and
nobler, where original disposition and outward lot have
worked together to such a degree that the criminal act can almost
be said to come as a necessary product of them, without
there having been any possibility for the individual's own laboring,
since consciousness has not been able to awaken at all. Such
an individual can be legally punishable and yet morally
% --- p. 332 ---
\opage{332}excusable: legally, one proceeds from the relation of the single individual
to society and society's right over against him; morally, the
gaze is directed solely at the relation between conscious striving and
the conditions. But the common conception sees the matter
in such a way that a criminal offense is a higher degree of a
merely moral transgression; it does not see that the properly ethical
moment can often be wholly subordinate in what is legally
punishable. Here, then, society can be within its right in
holding to the legal way of viewing; but the moral one
must not be overlooked either: it enjoins at once a reconciling
conception of the ``criminal'' and points to where
the real difference between good and evil comes forth.
The medical way of viewing is indisputably right over against
the legal one when it maintains ``that a considerable part of
criminals are feeble-minded or epileptic or come from
families in which there is mental illness, epilepsy, or some
other nervous disease''. But it often becomes guilty of the same
confusion as the vulgar way of viewing, only
in another manner. Because doubt can be raised about the moral
accountability of criminal offenses in certain cases, this
accountability is not therefore abolished altogether. As if
it found application only where there is legal talk of
crime!

Maudsley defines the physician's standpoint as the practical one,
where mind and body are regarded as a unity of essence;
only purely theoretically can one separate the mind from the body and
immerse oneself in self-consciousness. In practical respects we
perhaps also come to agree with him more than in theory.
It is beyond all doubt that the physician is called, and will ever
more and more be called, to a practical office of the
greatest humane significance. But precisely for that reason it is also
of the greatest significance that his cultivation and way of seeing
do not become one-sided. By the practical standpoint
Maudsley seems to understand the physiological in opposition to what
he calls the metaphysical. And yet the physiological
standpoint too is still only a theoretically one-sided one when it is not
% --- p. 333 ---
\opage{333}combined with an actual, living psychological insight; this
has been amply shown in the foregoing. The physician does indeed often
acquire such a psychological insight, since through close
attention to the states of the body he also sharpens
his eye for the observation of the mental; but penetrating
and complete it can only become when the psychological
is recognized as a circle of laws and conditions peculiar in itself,
which is not merely an imprint of the bodily states.
But things do not yet seem to have come that far. ``It
is rare,'' says Stuart Mill, ``to find a physician who is also
a psychologist.'' Even the founder of positivism, Auguste
Comte, shared this opinion; he regarded physicians only
as veterinarians, because they concern themselves solely with man's
animal, not with his human, nature.

The truly practical standpoint here too points
back to the true theory, in that it demands that all considerations be
respected. Therefore positivism and materialism,
which have grasped only one side of life and of the essence
of existence and are thus one-sided in theory, cannot lead to
the true practice either.

\begin{flushright}\textbf{Harald Høffding.}\end{flushright}
\begin{center}\rule{2.2cm}{0.4pt}\\[-10pt]\rule{2.2cm}{0.4pt}\end{center}

\end{document}
